% ============================================================
%  Setmana 7 — Unitat 8 · Derivades
%  Exercicis: 46, 54, 58, 93, 97
% ============================================================

\setmana{7}{Unitat 8 \textperiodcentered{} Derivades}{1 de novembre de 2026}%
        {46, 54, 58, 93 i 97}

\begin{errada}{Exercici 58}
  \ondiu{«Els punts de tall de la funció amb els eixos són:»}
  \hadedir{«Els punts de tall de la recta tangent amb els eixos són:»}
  \motiu{Els punts que es calculen a continuació, $(0,7)$ i
    $\left(\tfrac{7}{2},0\right)$, són els talls de la recta tangent
    $y=-2x+7$, que és la que forma el triangle amb els eixos. La funció $f$
    talla els eixos en uns altres punts: $(0,-1)$ i $(-1,0)$.}
\end{errada}

%\senseerrades{46, 54, 93 i 97}

% ---------- Observacions (només surten amb \obstrue) ----------
\ifobs
\bigskip
\noindent
Dues observacions menors:\par\nopagebreak

\begin{observacio}{Exercici 54}
  L'enunciat escriu la recta com a $f(x)=2x-123$, amb la mateixa lletra $f$
  de la funció $f(x)=x^3-x+4$. Són dues coses diferents: la recta és
  $y=2x-123$, que és com l'escriu la solució.
\end{observacio}

\begin{observacio}{Exercici 58}
  A la solució hi ha una errada tipogràfica sense importància: hi diu
  «traingle» en lloc de «triangle».
\end{observacio}
\fi
